% =====================================================================
\chapter{La base de données}
\label{ch:db}

\section{Organisation générale}

Le schéma comporte \textbf{16 tables} réparties en cinq familles, définies par
16 classes réparties sur 11 fichiers de modèle. Toutes héritent de
\fn{Base = declarative\_base()} (\file{app/database/base.py}) et sont importées
dans \file{app/models/\_\_init\_\_.py}, ce qui garantit qu'Alembic les voit
toutes.

\begin{longtable}{@{}p{3.4cm}p{11.4cm}@{}}
\toprule
\textbf{Famille} & \textbf{Tables}\\
\midrule
\endhead
Identité \& sécurité & \tbl{roles}, \tbl{users}, \tbl{audit\_logs}, \tbl{access\_requests}\\
Incidents \& actions & \tbl{incident\_tickets}, \tbl{incident\_ticket\_notes}, \tbl{action\_requests}, \tbl{mcp\_calls}\\
Conversation & \tbl{chat\_conversations}, \tbl{chat\_messages}\\
Demandes de service & \tbl{service\_requests}, \tbl{request\_messages}, \tbl{request\_events}, \tbl{request\_attachments}\\
Connaissance \& mesures & \tbl{knowledge\_documents}, \tbl{disk\_metrics}\\
\bottomrule
\end{longtable}

\begin{infobox}
\textbf{Deux conventions de clé primaire coexistent.} Les tables héritées de la
V1 (\tbl{users}, \tbl{roles}, \tbl{audit\_logs}, \tbl{access\_requests})
utilisent des entiers auto-incrémentés ; les tables métier ajoutées ensuite
utilisent des \texttt{UUID}. Les commentaires du code portent la trace de cette
transition : plusieurs modèles annoncent « \textsc{corrigé} : Integer (users.id
est un Integer, pas un UUID) », après un premier jet où les clés étrangères
vers \tbl{users} avaient été déclarées en UUID --- ce que PostgreSQL refuse.

Enfin, \tbl{chat\_conversations} et \tbl{chat\_messages} utilisent des
\texttt{BigInteger} auto-incrémentés, et ce choix est justifié : c'est
l'identifiant qui donne \emph{l'ordre du fil}, parce que les deux messages d'un
même tour partagent le \texttt{created\_at} renvoyé par \texttt{now()}.
\end{infobox}

\section{Schéma relationnel}

\begin{figure}[h!]
\centering
\begin{tikzpicture}[node distance=9mm and 13mm, every node/.style={font=\scriptsize\sffamily}]

\node[bluebox,minimum width=24mm] (roles) {\textbf{roles}\\\tiny id · name};
\node[bluebox,below=of roles,minimum width=24mm] (users) {\textbf{users}\\\tiny id · email · hashed\_password\\\tiny full\_name · is\_active · role\_id};

\node[redbox,right=20mm of users,minimum width=30mm] (tickets) {\textbf{incident\_tickets}\\\tiny id(UUID) · ticket\_number\\\tiny severity · status · mcp\_data(JSONB)\\\tiny recommendations(JSONB) · created\_by};
\node[redbox,below=of tickets,minimum width=30mm] (notes) {\textbf{incident\_ticket\_notes}\\\tiny id · ticket\_id · author\_id · content};
\node[ambox,above=of tickets,minimum width=30mm] (actions) {\textbf{action\_requests}\\\tiny id(UUID) · action\_type · target\\\tiny parameters(JSONB) · status · result(JSONB)\\\tiny requested\_by · confirmed\_by\\\tiny incident\_ticket\_id · chat\_message\_id};
\node[ambox,above=of actions,minimum width=30mm] (mcpcalls) {\textbf{mcp\_calls}\\\tiny id(UUID) · persona · tool\_name\\\tiny arguments(JSONB) · result(JSONB)\\\tiny status · correlation\_id\\\tiny user\_id · action\_request\_id};

\node[greenbox,right=20mm of tickets,minimum width=28mm] (conv) {\textbf{chat\_conversations}\\\tiny id(BigInt) · user\_id · title\\\tiny incident\_ticket\_id};
\node[greenbox,above=of conv,minimum width=28mm] (msg) {\textbf{chat\_messages}\\\tiny id(BigInt) · conversation\_id\\\tiny role · content};

\node[vibox,below=of notes,minimum width=28mm] (sreq) {\textbf{service\_requests}\\\tiny id(UUID) · direction · status\\\tiny created\_by · assigned\_to};
\node[vibox,right=20mm of sreq,minimum width=26mm] (satt) {\textbf{request\_attachments}\\\tiny id · request\_id · stored\_path};

\draw[flow] (roles)--(users);
\draw[flow] (users)--node[lbl,above]{created\_by}(tickets);
\draw[flow] (tickets)--(notes);
\draw[flow] (tickets)--(actions);
\draw[flow] (actions)--(mcpcalls);
\draw[flow] (tickets)--node[lbl,above]{investigation}(conv);
\draw[flow] (conv)--(msg);
\draw[flow] (msg.west) to[bend right=20] node[lbl,left]{chat\_message\_id} (actions.east);
\draw[flow] (users.south) to[bend right=25] (sreq.west);
\draw[flow] (sreq)--(satt);
\end{tikzpicture}
\caption{Relations principales. \tbl{audit\_logs}, \tbl{access\_requests},
\tbl{knowledge\_documents} et \tbl{disk\_metrics} ne sont pas représentées :
elles n'ont au plus qu'une clé étrangère vers \tbl{users}.}
\end{figure}

\clearpage

\section{Table par table}

% -------------------------------------------------- roles
\subsection{\tbl{roles} --- \file{models/role.py}}

\textbf{Rôle} : liste fermée des rôles applicatifs.\\
\textbf{Clé primaire} : \texttt{id} (Integer).\\
\textbf{Colonnes} : \texttt{id}, \texttt{name} (String(50), \texttt{unique},
\texttt{NOT NULL}).\\
\textbf{Relations} : \texttt{users} (un-à-plusieurs, \texttt{back\_populates}).\\
\textbf{Données} : trois lignes créées par \file{database/seed.py} ---
\texttt{linux\_engineer}, \texttt{storage\_engineer}, \texttt{administrator}.\\
\textbf{Utilisée dans} : \fn{require\_role}, \fn{request\_service.\_users\_with\_role},
\fn{admin\_service.list\_roles}, \fn{incident\_watcher.\_system\_user\_id}.

% -------------------------------------------------- users
\subsection{\tbl{users} --- \file{models/user.py}}

\textbf{Rôle} : comptes applicatifs.\\
\textbf{Clé primaire} : \texttt{id} (Integer, indexé).\\
\textbf{Clé étrangère} : \texttt{role\_id} \arrowto{} \texttt{roles.id}
(\texttt{NOT NULL}).\\
\textbf{Colonnes} : \texttt{email} (unique, indexé), \texttt{hashed\_password}
(bcrypt), \texttt{full\_name}, \texttt{is\_active} (défaut \texttt{True}),
\texttt{created\_at}.\\
\textbf{Utilisée dans} : \fn{authenticate\_user}, \fn{get\_current\_user},
\fn{\_extract\_user\_id} du middleware, et comme cible de \emph{toutes} les clés
étrangères d'auteur.

\begin{infobox}
Le mot de passe n'est jamais stocké en clair : \fn{hash\_password} utilise
\texttt{CryptContext(schemes=["bcrypt"])}. Le compte de service
\texttt{oasis-system@oasis.local} reçoit
\texttt{hash\_password(secrets.token\_urlsafe(32))} et
\texttt{is\_active=False}.
\end{infobox}

% -------------------------------------------------- audit_logs
\subsection{\tbl{audit\_logs} --- \file{models/audit\_log.py}}

\textbf{Rôle} : journal immuable des requêtes HTTP.\\
\textbf{Clé primaire} : \texttt{id} (Integer).\\
\textbf{Clé étrangère} : \texttt{user\_id} \arrowto{} \texttt{users.id}
(\emph{nullable} --- une requête anonyme est aussi journalisée).\\
\textbf{Colonnes} : \texttt{action} (méthode HTTP), \texttt{resource} (chemin),
\texttt{details} (JSON : \texttt{status\_code}, \texttt{duration\_ms},
\texttt{client\_ip}), \texttt{timestamp}.\\
\textbf{Écrite par} : \cls{AuditMiddleware} uniquement.\\
\textbf{Lue par} : \fn{audit\_service.list\_recent} \arrowto{}
\route{GET /admin/audit-logs}.

\begin{keybox}
Le docstring du modèle pose une règle : « cette table est pensée pour être
\emph{insert-only} --- on n'update et on ne supprime JAMAIS une ligne d'audit
log, on ne fait qu'en ajouter ». Aucun code du dépôt ne la modifie ni ne la
supprime : la règle est tenue.
\end{keybox}

% -------------------------------------------------- access_requests
\subsection{\tbl{access\_requests} --- \file{models/access\_request.py}}

\textbf{Rôle} : demande d'accès temporaire en lecture seule à un autre
workspace.\\
\textbf{Clé primaire} : \texttt{id} (Integer).\\
\textbf{Clés étrangères} : \texttt{requested\_by\_id} et \texttt{reviewed\_by\_id}
\arrowto{} \texttt{users.id} (deux relations distinctes, désambiguïsées par
\texttt{foreign\_keys=[...]}).\\
\textbf{Colonnes} : \texttt{target\_workspace} (texte libre), \texttt{reason},
\texttt{status} (\texttt{pending} / \texttt{approved} / \texttt{rejected} /
\texttt{expired}), \texttt{granted\_token}, \texttt{expires\_at},
\texttt{created\_at}, \texttt{reviewed\_at}.\\
\textbf{Utilisée dans} : \file{services/access\_request\_service.py} et
\route{/cross-domain}.

\begin{warnbox}
Le jeton généré à l'approbation est stocké dans \texttt{granted\_token}, mais
\textbf{aucun code du dépôt ne le vérifie ensuite}. Le statut \texttt{EXPIRED} de
l'énumération n'est jamais affecté non plus. Ce workflow est complet côté
enregistrement, incomplet côté application des droits.
\end{warnbox}

% -------------------------------------------------- incident_tickets
\subsection{\tbl{incident\_tickets} --- \file{models/incident\_ticket.py}}

\textbf{Rôle} : le ticket d'incident, pivot du cycle de vie.\\
\textbf{Clé primaire} : \texttt{id} (UUID).\\
\textbf{Clé étrangère} : \texttt{created\_by} \arrowto{} \texttt{users.id}
(\texttt{NOT NULL}).

\begin{longtable}{@{}p{3.6cm}p{2.6cm}p{8.4cm}@{}}
\toprule
\textbf{Colonne} & \textbf{Type} & \textbf{Contenu}\\
\midrule
\endhead
\texttt{ticket\_number} & String unique & \texttt{INC-2026-0001}, calculé par \fn{\_generate\_ticket\_number}\\
\texttt{title} & String & Ex. « Backup job backup-app-critical.timer failed »\\
\texttt{description} & Text & Phrase générée par la règle de détection\\
\texttt{severity} & Enum & \texttt{critical} / \texttt{high} / \texttt{medium} / \texttt{low}\\
\texttt{status} & Enum & \texttt{open} / \texttt{in\_progress} / \texttt{resolved} / \texttt{closed}\\
\texttt{affected\_system} & String & \textbf{Libellé lisible} (Description systemd)\\
\texttt{root\_cause} & Text & Rempli par la synthèse post-investigation\\
\texttt{impact\_summary} & Text & Idem\\
\texttt{mcp\_data} & \textbf{JSONB} & Relevé brut de l'outil au moment de la détection --- porte \texttt{job\_id} ou \texttt{volume\_id}, c'est-à-dire l'\textbf{identifiant technique}\\
\texttt{recommendations} & \textbf{JSONB} & Liste d'objets \texttt{\{order, description, action\_type, target, requires\_confirmation\}}\\
\texttt{correlation\_id} & UUID & Relie le ticket aux appels MCP de la même analyse\\
\texttt{created\_at} / \texttt{updated\_at} / \texttt{resolved\_at} & Timestamptz & Horodatages\\
\bottomrule
\end{longtable}

\textbf{Relations} : \texttt{creator} (\cls{User}), \texttt{notes} (cascade
\texttt{all, delete-orphan}), \texttt{action\_requests}.

\begin{keybox}
\textbf{La distinction à retenir absolument.} \texttt{affected\_system} est un
\emph{libellé} (« Sauvegarde de l'application critique »),
\texttt{mcp\_data["job\_id"]} est l'\emph{identifiant technique}
(\texttt{backup-app-critical.timer}). C'est le second que les outils attendent.
Toute la logique de \fn{IncidentContext.technical\_id} et de correction de cible
dans \fn{\_propose\_action} existe pour ne pas les confondre.
\end{keybox}

% -------------------------------------------------- notes
\subsection{\tbl{incident\_ticket\_notes}}

\textbf{Rôle} : piste d'audit lisible du ticket.\\
\textbf{Clé primaire} : \texttt{id} (UUID).\\
\textbf{Clés étrangères} : \texttt{ticket\_id} \arrowto{}
\texttt{incident\_tickets.id} ; \texttt{author\_id} \arrowto{} \texttt{users.id}.\\
\textbf{Écrite à} : détection automatique, ouverture d'investigation,
confirmation, rejet, échec d'action, clôture.\\
\textbf{Pourquoi} : « c'est ce qui permet de relire après coup pourquoi un
incident a été clos, et par qui --- l'horodatage seul des colonnes du ticket ne
le dit pas ».

% -------------------------------------------------- action_requests
\subsection{\tbl{action\_requests} --- \file{models/action\_request.py}}

\textbf{Rôle} : \textbf{la table qui matérialise la confirmation humaine.}\\
\textbf{Clé primaire} : \texttt{id} (UUID).

\begin{longtable}{@{}p{3.8cm}p{2.9cm}p{7.9cm}@{}}
\toprule
\textbf{Colonne} & \textbf{Type} & \textbf{Contenu}\\
\midrule
\endhead
\texttt{action\_type} & String indexé & Nom exact de l'outil MCP\\
\texttt{target} & String & Cible technique\\
\texttt{parameters} & \textbf{JSONB} & Arguments, \emph{sans} \texttt{confirm}\\
\texttt{status} & Enum & \texttt{pending} / \texttt{confirmed} / \texttt{completed} / \texttt{failed} / \texttt{rejected}\\
\texttt{requested\_by} & FK users & Auteur de la proposition\\
\texttt{confirmed\_by} & FK users, nullable & Qui a cliqué « Confirm »\\
\texttt{incident\_ticket\_id} & FK, nullable & Ticket rattaché\\
\texttt{chat\_message\_id} & FK BigInteger, \texttt{ON DELETE SET NULL} & Message qui porte la proposition\\
\texttt{result} & \textbf{JSONB} & Réponse brute de l'outil, ou \texttt{\{"error": ...\}}\\
\texttt{created\_at} / \texttt{confirmed\_at} / \texttt{completed\_at} & Timestamptz & Les trois temps de l'action\\
\bottomrule
\end{longtable}

\textbf{Pourquoi \texttt{chat\_message\_id}} : le commentaire du modèle est
explicite --- « Sert à réafficher les boutons Confirmer/Rejeter quand on rouvre
le fil : sans lui, une proposition disparaîtrait de l'écran au rechargement
alors qu'elle reste en attente en base. »

% -------------------------------------------------- mcp_calls
\subsection{\tbl{mcp\_calls} --- \file{models/mcp\_call.py}}

\textbf{Rôle} : preuve qu'une réponse de l'agent s'appuie sur un relevé réel.\\
\textbf{Clé primaire} : \texttt{id} (UUID).\\
\textbf{Clés étrangères} : \texttt{user\_id} \arrowto{} \texttt{users.id}
(\texttt{NOT NULL}, indexé) ; \texttt{action\_request\_id} \arrowto{}
\texttt{action\_requests.id} (nullable).\\
\textbf{Colonnes} : \texttt{persona} (« storage »), \texttt{tool\_name},
\texttt{arguments} (JSONB), \texttt{result} (JSONB), \texttt{status}
(\texttt{success} / \texttt{failed} / \texttt{denied}), \texttt{error\_message},
\texttt{correlation\_id} (indexé), \texttt{created\_at}.\\
\textbf{Écrite par} : \fn{mcp\_audit.log\_mcp\_call}, appelée \emph{uniquement}
depuis \fn{StorageMCPClient.call\_tool}.

\begin{infobox}
Le commentaire de \texttt{arguments} précise : ce sont « les arguments réellement
envoyés au tool (après normalisation et garde-fou, donc reflète le
\texttt{confirm=True/False} réel, pas ce que le LLM avait proposé) ». C'est
exactement ce qu'il faut pour un audit de sécurité.
\end{infobox}

% -------------------------------------------------- chat
\subsection{\tbl{chat\_conversations} et \tbl{chat\_messages}}

\textbf{\tbl{chat\_conversations}} --- clé primaire \texttt{id} (BigInteger) ;
\texttt{user\_id} \arrowto{} \texttt{users.id} ; \texttt{title}
(\(\le\) 120 caractères, jamais vide) ; \texttt{incident\_ticket\_id}
(\texttt{ON DELETE SET NULL}) ; \texttt{created\_at} ; \texttt{updated\_at}.

\textbf{\tbl{chat\_messages}} --- clé primaire \texttt{id} (BigInteger) ;
\texttt{conversation\_id} \arrowto{} \texttt{chat\_conversations.id}
(\texttt{ON DELETE CASCADE}) ; \texttt{role} (\texttt{"user"} ou
\texttt{"assistant"}, stocké en texte) ; \texttt{content} ; \texttt{created\_at}.

\begin{keybox}
\textbf{Index unique partiel.} La migration \file{b6d0f3a91c47} crée
\texttt{uq\_chat\_conversations\_incident\_ticket} sur
\texttt{(incident\_ticket\_id, user\_id)}, unique, avec
\texttt{postgresql\_where = incident\_ticket\_id IS NOT NULL}. Justification :
« Un ingénieur n'a qu'un fil d'investigation par ticket : l'unicité est ce qui
permet de rouvrir l'existant sans arbitrer entre plusieurs candidats \ldots{}
L'index est partiel, la colonne étant nulle pour les fils ordinaires, qui sont
la majorité. »
\end{keybox}

\noindent
\textbf{Pourquoi le propriétaire n'est pas dupliqué sur le message} : « Le
propriétaire se lit sur la conversation : le dupliquer ici ouvrirait la porte à
un message rattaché à un fil qui n'est pas le sien. »

% -------------------------------------------------- service_requests
\subsection{\tbl{service\_requests} et ses satellites}

\textbf{\tbl{service\_requests}} --- clé primaire UUID.
Clés étrangères : \texttt{created\_by} et \texttt{assigned\_to} \arrowto{}
\texttt{users.id} (les deux \texttt{NOT NULL} et indexées).
Colonnes notables : \texttt{direction} (Enum), \texttt{title},
\texttt{description} (nullable), \texttt{request\_type}, \texttt{priority},
\texttt{status}, \texttt{deadline}, \texttt{administrator\_response},
\texttt{decided\_at}, \texttt{acknowledged\_at}, \texttt{completed\_at}.

\begin{longtable}{@{}p{4.0cm}p{10.8cm}@{}}
\toprule
\textbf{Table satellite} & \textbf{Statut réel}\\
\midrule
\endhead
\tbl{request\_events} & \textbf{Active}. Une ligne par transition (\texttt{from\_status}, \texttt{to\_status}, \texttt{note}, \texttt{actor\_id}). Jamais renvoyée au client : « l'interface se contente du statut courant et de son horodatage ».\\
\tbl{request\_attachments} & \textbf{Active}. Fichier écrit sous l'UUID de la ligne ; le nom d'origine ne survit que comme métadonnée, « ce qui rend toute traversée de répertoire impossible ».\\
\tbl{request\_messages} & \textbf{Gelée}. Le docstring l'assume : « le fil a été retiré de l'interface et de l'API. La table est conservée pour ne pas détruire les échanges déjà enregistrés, mais plus rien n'y écrit. »\\
\bottomrule
\end{longtable}

\begin{infobox}
\textbf{Trois modèles au nom proche, à ne pas confondre --- le code le dit
lui-même.} \cls{AccessRequest} = accès temporaire cross-domain ;
\cls{ActionRequest} = exécution d'un outil MCP sensible ;
\cls{ServiceRequest} = canal de communication humain, qui « ne déclenche aucune
action technique ».
\end{infobox}

% -------------------------------------------------- knowledge / metrics
\subsection{\tbl{knowledge\_documents} --- \file{models/knowledge\_document.py}}

\textbf{Rôle} : métadonnées des documents du Data Hub. Le fichier lui-même vit
sur disque ; « cette table ne garde que les métadonnées et l'état d'indexation
Chroma ».\\
\textbf{Clé primaire} : \texttt{id} (UUID) --- c'est aussi le nom du fichier sur
disque.\\
\textbf{Clé étrangère} : \texttt{uploaded\_by} \arrowto{} \texttt{users.id}.\\
\textbf{Colonnes} : \texttt{filename}, \texttt{stored\_path},
\texttt{content\_type}, \texttt{size\_bytes}, \texttt{collection} (défaut
\texttt{storage\_kb}, indexé), \texttt{chunk\_count}, \texttt{status}
(\texttt{pending} / \texttt{indexed} / \texttt{failed}), \texttt{error},
\texttt{created\_at}, \texttt{indexed\_at}.

\subsection{\tbl{disk\_metrics} --- \file{models/metric.py}}

\textbf{Rôle} : un échantillon de performance, pour un device, à un instant.\\
\textbf{Clé primaire} : \texttt{id} (UUID).\\
\textbf{Aucune clé étrangère.}\\
\textbf{Colonnes} : \texttt{device} (indexé), \texttt{iops},
\texttt{throughput\_mbps}, \texttt{latency\_ms}, \texttt{recorded\_at} (indexé).\\
\textbf{Écrite par} : \fn{metrics\_service.collect\_once}.
\textbf{Lue par} : \fn{read\_series}. \textbf{Purgée par} : \fn{purge\_old}
(rétention 31 jours).\\
\textbf{Principe} : « l'historique 24h/7j/30j se construit au fil des
prélèvements, il n'est jamais fabriqué ».

\section{Les migrations Alembic}

Quatorze migrations forment une \textbf{chaîne linéaire}, sans branche :

\begin{lstlisting}[caption={Ordre réel des révisions}]
a5bd2724f8a3  create users and roles tables            (base, down_revision = None)
a7678a14436e  add audit_logs table
8aab52e3789f  add access_requests table
cfa440d22089  add incidents table            <-- MIGRATION VIDE (upgrade: pass)
c93419b63533  add incident_tickets and action_requests tables
d41f7c2b9e10  add knowledge_documents table (Data Hub)
e58a3d9c4712  add mcp_calls table
f7c1b0d5a983  add disk_metrics table (real performance history)
b2f4d6c81a05  add chat_messages table (historique de conversation)
c7e1a94b3d20  add chat_conversations (plusieurs fils par utilisateur)
d9a3f17c6b42  add service_requests, request_messages, request_events
f4b28c50e173  add decided_at to service_requests
a1c7e94d2f80  add request_attachments
b6d0f3a91c47  link investigation chat to incidents            <-- HEAD
\end{lstlisting}

\begin{warnbox}
\textbf{Détail à ne pas cacher.} La révision \file{cfa440d22089}, malgré son
titre « add incidents table », a un \fn{upgrade()} qui ne contient qu'un
\texttt{pass} : c'est un squelette généré puis abandonné. Les vraies tables
d'incident arrivent avec \file{c93419b63533}.
\end{warnbox}

\section{Workflow base de données : le cycle complet d'une action}

Reprenons la demande de l'énoncé --- \emph{création \arrowto{} stockage
\arrowto{} récupération \arrowto{} modification du statut \arrowto{} affichage}
--- appliquée à une \cls{ActionRequest} réelle.

\begin{longtable}{@{}p{0.7cm}p{4.4cm}p{9.6cm}@{}}
\toprule
\textbf{\#} & \textbf{Étape} & \textbf{Effet exact en base}\\
\midrule
\endhead
1 & Le LLM appelle \fn{run\_backup} & Rien encore : l'appel est intercepté par \fn{\_propose\_action}\\
2 & \fn{create\_action\_request} & \texttt{INSERT} dans \tbl{action\_requests} : \texttt{status='pending'}, \texttt{requested\_by}, \texttt{target}, \texttt{parameters}, \texttt{incident\_ticket\_id}. \texttt{commit} + \texttt{refresh}\\
3 & \fn{chat\_service.save\_turn} & \texttt{INSERT} de deux \tbl{chat\_messages} dans la même transaction ; \texttt{chat\_conversations.updated\_at = now()}\\
4 & \fn{attach\_actions} & \texttt{UPDATE action\_requests SET chat\_message\_id = <id du message assistant> WHERE id IN (...)}\\
5 & Rechargement du fil & \texttt{SELECT} des messages avec leur relation \texttt{action\_requests} \arrowto{} \cls{ChatMessageOut} \arrowto{} boutons rendus par \cls{ActionProposal}\\
6 & Clic « Confirm » & \texttt{UPDATE} : \texttt{status='confirmed'}, \texttt{confirmed\_by}, \texttt{confirmed\_at}. \textbf{commit immédiat}\\
7 & Appel MCP & \texttt{INSERT} dans \tbl{mcp\_calls} : \texttt{tool\_name}, \texttt{arguments} (avec \texttt{confirm=True}), \texttt{result}, \texttt{status}, \texttt{action\_request\_id}\\
8 & Qualification & \texttt{UPDATE} : \texttt{status='completed'} ou \texttt{'failed'}, \texttt{result} (JSONB), \texttt{completed\_at}\\
9 & \fn{\_close\_the\_loop} & \texttt{INSERT} dans \tbl{incident\_ticket\_notes} avec le verdict\\
10 & Si résolu & \texttt{UPDATE incident\_tickets} : \texttt{status='resolved'}, \texttt{resolved\_at} ; plus une note de clôture\\
11 & Sinon, plus tard & \fn{incident\_watcher.\_verify\_resolved} refait le test au passage suivant (300\,s) et clôt si l'indicateur est redevenu sain\\
\bottomrule
\end{longtable}

\section{Workflow base de données : naissance d'un incident}

\begin{enumerate}
  \item \fn{list\_backups()} lit les timers systemd par SSH.
  \item \fn{check\_backup\_status(job)} renvoie un dictionnaire de ticket, ou
  \texttt{None}.
  \item \fn{\_latest\_ticket} + \fn{\_already\_handled} : \texttt{SELECT} du
  dernier ticket sur ce système, pour ne pas rouvrir un incident déjà traité.
  \item \fn{create\_ticket\_from\_diagnosis} \arrowto{} \fn{\_open\_ticket\_for}
  fait un \texttt{SELECT} sur \texttt{mcp\_data['job\_id']} (opérateur JSONB
  \texttt{->>}) parmi les statuts \texttt{open} et \texttt{in\_progress}.
  \begin{itemize}
    \item \textbf{Trouvé} \arrowto{} \texttt{UPDATE} du titre, de la description,
    de la sévérité, de \texttt{mcp\_data} et des recommandations. Un seul ticket
    pour une anomalie qui persiste.
    \item \textbf{Absent} \arrowto{} \fn{\_generate\_ticket\_number}
    (\texttt{SELECT max(substr(ticket\_number, 10))}) puis \texttt{INSERT}.
  \end{itemize}
  \item Si le ticket est neuf : \texttt{INSERT} d'une première note d'audit.
\end{enumerate}
